\section{Conclusion}\label{sec:conclusion}
We presented a closed-loop quality-triage system for skin-lesion diagnosis on degraded
consumer imagery, built on two contributions. The first is a feature-level
diagnosis-preserving enhancement objective (DP-Loss) whose information loss is bounded by a
mutual-information lower bound (Lemma~\ref{lem:dp}) and confirmed empirically. The second is
a query-for-retake channel that, to our knowledge, is the first to request re-acquisition
rather than merely abstain or defer, and it closes the acquisition loop under a local,
per-band conditional risk guarantee (Theorem~\ref{thm:agent-compact}). We also chart its
limits. Enhancement turns net-negative on extreme contrast loss and on malignant lesions
under severe degradation, and global triage does not yet beat direct diagnosis. We read
these honest negatives as motivation for learned routing thresholds. Mapping where
degradation-aware enhancement helps, and where it must yield to re-acquisition, is the more
useful contribution for deploying diagnosis on the quality-varying inputs that curated
benchmarks exclude by design.
